\documentclass[11pt]{article}
\usepackage[margin=1in]{geometry}
\usepackage{amsmath}
\usepackage{amssymb}
\usepackage{enumitem}
\usepackage{fancyhdr}
\usepackage{graphicx}
\usepackage{tikz}

\pagestyle{fancy}
\fancyhf{}
\lhead{AP Statistics - Unit 2.9}
\rhead{Analyzing Departures from Linearity}
\cfoot{\thepage}

\begin{document}

\begin{center}
{\Large \textbf{Follow-Along Worksheet: Topic 2.9}}\\[0.5em]
{\large Analyzing Departures from Linearity}\\[1em]
Name: \underline{\hspace{3in}} \quad Period: \underline{\hspace{1in}}
\end{center}

\section*{Part 1: Influential Points (Video 1)}

\subsection*{Opening (0:00--0:29)}
\begin{enumerate}[label=\arabic*.]
\item What three main concepts will this video cover?
\begin{enumerate}[label=\alph*)]
\item \underline{\hspace{3.5in}}
\item \underline{\hspace{3.5in}}
\item \underline{\hspace{3.5in}}
\end{enumerate}
\end{enumerate}

\subsection*{Context and Setup (0:30--1:41)}
\begin{enumerate}[resume, label=\arabic*.]
\item What relationship is being explored with the San Antonio grocery store data?

\vspace{0.5em}
$x$-variable: \underline{\hspace{3in}}

$y$-variable: \underline{\hspace{3in}}

\item Record the regression equation parameters:
\begin{itemize}
\item $y$-intercept $\approx$ \underline{\hspace{0.5in}}
\item slope $\approx$ \underline{\hspace{0.5in}}
\end{itemize}

\item The mean income is \$\underline{\hspace{1in}} and the mean number of items is \underline{\hspace{0.5in}}.

\textbf{Key Insight:} This point $(\bar{x}, \bar{y})$ always lies \underline{\hspace{2in}}.
\end{enumerate}

\subsection*{Low vs. High Leverage Points (1:42--3:41)}
\begin{enumerate}[resume, label=\arabic*.]
\item \textbf{Low leverage points} are close to \underline{\hspace{1.5in}} on the $x$-axis.

When removed, what happens to the regression line? \underline{\hspace{2.5in}}

\item \textbf{High leverage points} are far from \underline{\hspace{1.5in}} on the $x$-axis.

\item When three high leverage points were removed:
\begin{itemize}
\item The $y$-intercept changed from approximately \underline{\hspace{0.5in}} to \underline{\hspace{0.5in}}
\item This is a shift of \underline{\hspace{0.5in}} units!
\end{itemize}

\item \textbf{Definition:} High leverage points have unusually \underline{\hspace{1in}} or \underline{\hspace{1in}} $x$-values.
\end{enumerate}

\subsection*{Outliers in Regression (3:43--4:58)}
\begin{enumerate}[resume, label=\arabic*.]
\item Looking at the green point in the simple example:
\begin{itemize}
\item Is it a high leverage point? \underline{\hspace{0.5in}} 
\item Why or why not? \underline{\hspace{3in}}
\end{itemize}

\item An \textbf{outlier} in regression has an unusually high magnitude \underline{\hspace{1.5in}}.

\item When the outlier was removed, what happened to:
\begin{itemize}
\item Correlation ($r$): \underline{\hspace{1.5in}}
\item Coefficient of determination ($R^2$): \underline{\hspace{1.5in}}
\end{itemize}
\end{enumerate}

\subsection*{Types of Influential Points (5:02--6:03)}
\begin{enumerate}[resume, label=\arabic*.]
\item Complete the table for the three types of influential points:

\begin{center}
\begin{tabular}{|p{2in}|p{3.5in}|}
\hline
\textbf{Type} & \textbf{Primary Effect When Removed} \\
\hline
Outlier & Changes \underline{\hspace{2.5in}} \\
\hline
High Leverage Point & Changes \underline{\hspace{2.5in}} \\
\hline
Both Outlier AND High Leverage & Changes \underline{\hspace{2.5in}} \\
\hline
\end{tabular}
\end{center}

\item \textbf{Synthesis:} An influential point is defined as any point that, if removed, substantially changes the \underline{\hspace{0.8in}}, \underline{\hspace{0.8in}}, and/or \underline{\hspace{0.8in}}.
\end{enumerate}

\subsection*{Reflection Questions}
\begin{enumerate}[resume, label=\arabic*.]
\item Why might a point with an $x$-value close to $\bar{x}$ still be influential?

\vspace{1.5em}

\item Sketch a scatter plot showing:
\begin{itemize}
\item 5 points following a linear trend
\item 1 outlier that is NOT a high leverage point
\item 1 high leverage point that IS following the trend
\end{itemize}

\begin{center}
\begin{tikzpicture}[scale=0.8]
\draw[->] (0,0) -- (8,0) node[right] {$x$};
\draw[->] (0,0) -- (0,6) node[above] {$y$};
\foreach \x in {1,...,7}
    \draw (\x,0.1) -- (\x,-0.1) node[below] {\tiny \x};
\foreach \y in {1,...,5}
    \draw (0.1,\y) -- (-0.1,\y) node[left] {\tiny \y};
\end{tikzpicture}
\end{center}
\end{enumerate}

\newpage

\section*{Part 2: Transforming Data for Linearity (Video 2)}

\subsection*{Opening and Context (0:00--1:14)}
\begin{enumerate}[label=\arabic*.]
\item What will we learn about nonlinear relationships?
\begin{enumerate}[label=\alph*)]
\item What to do when data has a \underline{\hspace{2in}}
\item Effects of \underline{\hspace{1.5in}} on datasets
\item How to assess if transforming improved \underline{\hspace{1.5in}}
\end{enumerate}

\item The example explores the relationship between:
\begin{itemize}
\item $x$: \underline{\hspace{3in}} (GDP per person)
\item $y$: \underline{\hspace{3in}}
\end{itemize}
\end{enumerate}

\subsection*{Initial Analysis (1:15--3:06)}
\begin{enumerate}[resume, label=\arabic*.]
\item Describe the association: \underline{\hspace{1in}} (positive/negative/none)

\item Compare the United States and Japan:
\begin{center}
\begin{tabular}{|l|c|c|}
\hline
Country & Income per Person & Life Expectancy \\
\hline
United States & \$\underline{\hspace{0.8in}} & \underline{\hspace{0.5in}} years \\
\hline
Japan & \$\underline{\hspace{0.8in}} & \underline{\hspace{0.5in}} years \\
\hline
\end{tabular}
\end{center}

What does this comparison show? \underline{\hspace{3in}}

\item For the original linear model:
\begin{itemize}
\item $R^2 = $ \underline{\hspace{0.5in}}\%
\item The residual plot shows a \underline{\hspace{1.5in}} pattern (good/bad?)
\end{itemize}

\item What type of model does the speaker suggest might be more appropriate? \underline{\hspace{1.5in}}
\end{enumerate}

\subsection*{Data Transformation (3:07--4:11)}
\begin{enumerate}[resume, label=\arabic*.]
\item Why is income data typically right-skewed?

\vspace{1.5em}

\item What does a log transformation do to high values?
\begin{itemize}
\item Makes them \underline{\hspace{1.5in}}
\item While preserving the \underline{\hspace{1.5in}} between values
\end{itemize}

\item The transformation applied was: $x_{new} = $ \underline{\hspace{1in}}
\end{enumerate}

\subsection*{Comparing Models (4:12--5:14)}
\begin{enumerate}[resume, label=\arabic*.]
\item Complete the comparison table:

\begin{center}
\begin{tabular}{|l|c|c|}
\hline
\textbf{Measure} & \textbf{Untransformed} & \textbf{Log-Transformed} \\
\hline
Scatter plot form & Curved & \underline{\hspace{1.2in}} \\
\hline
Residual plot pattern & Shows pattern & \underline{\hspace{1.2in}} \\
\hline
$R^2$ value & \underline{\hspace{0.6in}}\% & \underline{\hspace{0.6in}}\% \\
\hline
\end{tabular}
\end{center}

\item \textbf{Key Question:} How do we know the transformation improved the model? List two pieces of evidence:
\begin{enumerate}[label=\alph*)]
\item \underline{\hspace{4in}}
\item \underline{\hspace{4in}}
\end{enumerate}
\end{enumerate}

\subsection*{Other Transformations (5:15--5:40)}
\begin{enumerate}[resume, label=\arabic*.]
\item Besides logarithm, what other transformations are mentioned?
\begin{itemize}
\item \underline{\hspace{1.5in}}
\item \underline{\hspace{1.5in}}
\end{itemize}

\item In AP Statistics, will you typically transform data yourself? \underline{\hspace{0.5in}}

Instead, you'll assess model fit using \underline{\hspace{1.5in}} and \underline{\hspace{0.5in}} values.
\end{enumerate}

\subsection*{Synthesis Questions}
\begin{enumerate}[resume, label=\arabic*.]
\item Why might a linear model be preferred even when data shows a curved pattern?

\vspace{1.5em}

\item A student claims: ``Since the $R^2$ went from 46.6\% to 71.1\%, the transformed model explains 71.1\% of the variation in life expectancy.'' 

Is this interpretation correct? Explain carefully.

\vspace{2em}

\item When would you choose to transform data? Circle all that apply:
\begin{enumerate}[label=\alph*)]
\item The residual plot shows a clear pattern
\item The $R^2$ value is below 50\%
\item The original relationship appears nonlinear
\item You want to use linear regression methods
\item The data contains outliers
\end{enumerate}
\end{enumerate}

\subsection*{Practice Problem}
\begin{enumerate}[resume, label=\arabic*.]
\item A researcher studying the relationship between city population ($x$) and number of coffee shops ($y$) finds:
\begin{itemize}
\item Original model: $\hat{y} = 12 + 0.003x$ with $R^2 = 0.42$
\item After log-transforming $x$: $\hat{y} = -85 + 42\log(x)$ with $R^2 = 0.78$
\end{itemize}

\begin{enumerate}[label=\alph*)]
\item Which model appears to be a better fit? Why?

\vspace{1.5em}

\item If a city has a population of 100,000, predict the number of coffee shops using the transformed model. (Use $\log(100000) \approx 5$)

\vspace{2em}
\end{enumerate}
\end{enumerate}

\section*{Summary: Key Concepts Checklist}

Check off each concept as you master it:

\begin{itemize}
\item[$\square$] I can identify high leverage points (unusual $x$-values)
\item[$\square$] I can identify outliers in regression (large residuals)
\item[$\square$] I can explain how influential points affect slope, $y$-intercept, and correlation
\item[$\square$] I understand when to consider transforming data
\item[$\square$] I can assess model fit using residual plots
\item[$\square$] I can assess model fit using $R^2$ values
\item[$\square$] I can interpret transformed regression models
\end{itemize}

\vspace{1em}
\textbf{Remember:} Be critical, be cautious, be compassionate, and avoid bad statistics!

\end{document}